\documentclass[11pt,a4paper]{article}

\usepackage[margin=1in]{geometry}
\usepackage{graphicx}
\usepackage{listings}
\usepackage{xcolor}
\usepackage{hyperref}
\usepackage{booktabs}
\usepackage{enumitem}
\usepackage{fancyhdr}
\usepackage{titlesec}
\usepackage{caption}
\usepackage{amsmath}

% Colors
\definecolor{codebg}{RGB}{245,245,245}
\definecolor{codeframe}{RGB}{200,200,200}
\definecolor{keyword}{RGB}{0,0,180}
\definecolor{comment}{RGB}{0,128,0}
\definecolor{string}{RGB}{163,21,21}

% Define YAML language for listings
\lstdefinelanguage{yaml}{
    keywords={true, false, null, yes, no},
    sensitive=false,
    comment=[l]{\#},
    morestring=[b]',
    morestring=[b]",
}

% Code listing style
\lstset{
    backgroundcolor=\color{codebg},
    frame=single,
    rulecolor=\color{codeframe},
    basicstyle=\ttfamily\small,
    keywordstyle=\color{keyword}\bfseries,
    commentstyle=\color{comment}\itshape,
    stringstyle=\color{string},
    breaklines=true,
    breakatwhitespace=true,
    tabsize=2,
    showstringspaces=false,
    numbers=left,
    numberstyle=\tiny\color{gray},
    numbersep=8pt,
    xleftmargin=15pt,
    framexleftmargin=15pt,
}

% Header/footer
\pagestyle{fancy}
\fancyhf{}
\fancyhead[L]{\small AVROS --- Nav2 Route Server Integration}
\fancyhead[R]{\small March 2026}
\fancyfoot[C]{\thepage}

% Section formatting
\titleformat{\section}{\Large\bfseries}{}{0em}{}[\titlerule]
\titleformat{\subsection}{\large\bfseries}{}{0em}{}

\hypersetup{
    colorlinks=true,
    linkcolor=blue!60!black,
    urlcolor=blue!60!black,
    citecolor=blue!60!black,
}

\begin{document}

% ===== TITLE =====
\begin{titlepage}
\centering
\vspace*{2cm}
{\Huge\bfseries Nav2 Route Server Integration\\[0.3cm] Debug \& Resolution Report\par}
\vspace{1.5cm}
{\Large AVROS --- Autonomous Vehicle ROS2 Platform\par}
\vspace{1cm}
{\large California State Polytechnic University, Pomona\par}
\vspace{2cm}
{\large\today\par}
\vspace{2cm}

\begin{tabular}{ll}
\toprule
\textbf{Platform} & NVIDIA Jetson Orin (JetPack 6) \\
\textbf{ROS Distribution} & ROS 2 Humble Hawksbill \\
\textbf{Nav2 Version} & 1.1.x (Humble) \\
\textbf{Route Server} & \texttt{nav2\_route} v1.1.20 \\
\textbf{BehaviorTree.CPP} & v3 (3.8.7) \\
\bottomrule
\end{tabular}

\vfill
\end{titlepage}

\tableofcontents
\newpage

% ===== SECTION 1: OVERVIEW =====
\section{Overview}

This document describes the integration of the Nav2 Route Server (\texttt{nav2\_route}) into the AVROS autonomous vehicle platform, the problems encountered during bench testing, and the fixes applied to achieve a working navigation pipeline.

The Route Server provides graph-based road navigation using a pre-built GeoJSON road graph of the Cal Poly Pomona campus. It replaces the previous \texttt{route\_planner\_node.py}, which required internet access at runtime, snapped the vehicle to the nearest intersection with no bridging waypoints, and made 200+ blocking service calls per route.

\subsection{Architecture}

The navigation stack consists of:
\begin{itemize}[noitemsep]
    \item \textbf{Route Server} (\texttt{nav2\_route}) --- graph-based road planner using GeoJSON campus network
    \item \textbf{Planner Server} (SmacPlannerHybrid) --- freespace Dubin-curve planner for first/last-mile
    \item \textbf{Controller Server} (Regulated Pure Pursuit) --- path-following controller
    \item \textbf{Behavior Tree} --- orchestrates planning and execution
    \item \textbf{Costmaps} --- local (VoxelLayer + LiDAR/camera) and global (ObstacleLayer + LiDAR)
    \item \textbf{Localization} --- EKF + NavSat transform with fixed datum
\end{itemize}

\subsection{Test Environment}

All testing was performed indoors on a bench with the vehicle stationary. Key constraints:
\begin{itemize}[noitemsep]
    \item No GPS fix --- robot position fixed at $(0, 0, 0)$ in map frame
    \item LiDAR detecting indoor walls and objects
    \item Camera not connected (D455 not plugged in)
    \item No actuator movement --- routing pipeline only
\end{itemize}

% ===== SECTION 2: PROBLEMS AND FIXES =====
\section{Problems Encountered and Resolutions}

\subsection{Problem 1: Path Contains 0 Poses}

\subsubsection*{Symptom}

The \texttt{ComputeRoute} action returned 1,445 valid poses when called directly via CLI, but the BT navigator's \texttt{FollowPath} node received a path with 0 poses.

\begin{lstlisting}[language=bash]
[bt_navigator]: Resulting plan has 0 poses in it
[controller_server]: Received a goal with 0 poses, ignoring
\end{lstlisting}

\subsubsection*{Root Cause}

\texttt{ComputeRoute} returns a path starting at the \textbf{nearest graph node} to the robot, not at the robot's current position. With the robot at $(0, 0)$ and the nearest graph node at $(-165, 71)$ --- 180 meters away --- the entire path fell outside the local costmap's 10m window.

The Regulated Pure Pursuit controller's \texttt{transformGlobalPlan()} method prunes all poses outside the local costmap extent before execution. Since no poses were within range, the controller received an empty path.

On ROS 2 Jazzy, dedicated BT nodes (\texttt{GetPoseFromPath}, \texttt{ConcatenatePaths}, \texttt{ArePosesNear}) exist to handle this ``first-mile'' gap. These nodes do \textbf{not} exist on Humble.

\subsubsection*{Fix}

Replaced \texttt{ComputeRoute} with \texttt{ComputePathToPose} in the Behavior Tree. \texttt{ComputePathToPose} always plans from the robot's current position, ensuring the path starts within the local costmap.

\begin{lstlisting}[language=XML,caption={navigate\_route\_graph.xml --- Before (broken)}]
<ComputeRoute goal="{goal}" path="{path}" use_poses="true"/>
\end{lstlisting}

\begin{lstlisting}[language=XML,caption={navigate\_route\_graph.xml --- After (working)}]
<ComputePathToPose goal="{goal}" path="{path}" planner_id="GridBased"/>
\end{lstlisting}

The Route Server remains active and available for future use when first-mile BT nodes are ported to Humble or the platform upgrades to Jazzy.

\textbf{Commit:} \texttt{8af6013}

% -----

\subsection{Problem 2: No Valid Path Found (Unknown Space)}

\subsubsection*{Symptom}

SmacPlannerHybrid failed to find any path to the goal:

\begin{lstlisting}[language=bash]
[planner_server]: GridBased: failed to create plan
  with tolerance 0.50, no valid path found
\end{lstlisting}

\subsubsection*{Root Cause}

The global costmap had \texttt{track\_unknown\_space: true}, which marks all unobserved cells as \texttt{NO\_INFORMATION} (cost 255). With no sensor data populating the costmap at startup, the entire planning area was unknown space. SmacPlannerHybrid treats unknown cells as impassable by default.

\subsubsection*{Fix}

Set \texttt{track\_unknown\_space: false} in the global costmap configuration. This treats all cells as free until sensor data proves otherwise --- the correct behavior for outdoor operation without a static map.

\begin{lstlisting}[language=yaml,caption={nav2\_params.yaml --- global\_costmap change}]
global_costmap:
  global_costmap:
    ros__parameters:
      track_unknown_space: false   # free until proven occupied
\end{lstlisting}

\textbf{Commit:} \texttt{27fab04}

% -----

\subsection{Problem 3: Transform Tolerance Too Tight}

\subsubsection*{Symptom}

The controller server failed with a TF lookup error:

\begin{lstlisting}[language=bash]
[controller_server]: Unable to transform robot pose
  into global plan's frame
\end{lstlisting}

\subsubsection*{Root Cause}

The \texttt{transform\_tolerance} parameter in the Regulated Pure Pursuit controller was set to 0.1 seconds. During bench testing with timing gaps between TF publishers, this was too strict.

\subsubsection*{Fix}

Increased \texttt{transform\_tolerance} to 1.0 seconds (the Nav2 default):

\begin{lstlisting}[language=yaml,caption={nav2\_params.yaml --- FollowPath controller}]
FollowPath:
  plugin: "nav2_regulated_pure_pursuit_controller::..."
  transform_tolerance: 1.0   # was 0.1
\end{lstlisting}

\textbf{Commit:} \texttt{c783af3}

% -----

\subsection{Problem 4: Route Server Lifecycle Race Condition}

\subsubsection*{Symptom}

The Route Server remained stuck in the \texttt{unconfigured} lifecycle state after launch. The lifecycle manager for the route server started but failed silently:

\begin{lstlisting}[language=bash]
[lifecycle_manager_route]: Configuring route_server
# ... no response, route_server stays unconfigured
\end{lstlisting}

Meanwhile, the main Nav2 lifecycle manager successfully activated all its nodes.

\subsubsection*{Root Cause}

The \texttt{lifecycle\_manager\_route} started at 0.12 seconds after launch --- before the \texttt{route\_server} had finished registering its lifecycle services with the DDS discovery system. The main Nav2 lifecycle manager worked because it naturally started later (1.6 seconds) due to parameter loading overhead.

\subsubsection*{Diagnosis}

Manual lifecycle management confirmed the route server itself was functional:

\begin{lstlisting}[language=bash]
# Manual configure + activate worked perfectly
ros2 lifecycle set /route_server configure
ros2 lifecycle set /route_server activate
# Route server loaded graph: 52 nodes, 113 edges
\end{lstlisting}

\subsubsection*{Fix}

Wrapped the lifecycle manager node in a \texttt{TimerAction} with a 5-second delay, giving the route server time to register its DDS services:

\begin{lstlisting}[language=python,caption={navigation.launch.py --- TimerAction delay}]
from launch.actions import TimerAction

route_lifecycle_manager = TimerAction(
    period=5.0,
    actions=[
        Node(
            package='nav2_lifecycle_manager',
            executable='lifecycle_manager',
            name='lifecycle_manager_route',
            parameters=[{
                'autostart': True,
                'node_names': ['route_server'],
                'bond_timeout': 10.0,
            }],
            output='screen',
        ),
    ],
)
\end{lstlisting}

\textbf{Commit:} \texttt{8e7797d}

% -----

\subsection{Problem 5: Missing BT Plugin List (Humble)}

\subsubsection*{Symptom}

The BT navigator failed to load certain BT node plugins at runtime.

\subsubsection*{Root Cause}

On ROS 2 Jazzy, BT node plugins are auto-discovered. On Humble, they must be explicitly listed under the \texttt{plugin\_lib\_names} parameter in the \texttt{bt\_navigator} configuration. The original config did not include this list.

\subsubsection*{Fix}

Added the full explicit plugin list (38 plugins) to \texttt{nav2\_params.yaml}:

\begin{lstlisting}[language=yaml,caption={nav2\_params.yaml --- bt\_navigator plugin list (excerpt)}]
bt_navigator:
  ros__parameters:
    plugin_lib_names:
      - nav2_compute_path_to_pose_action_bt_node
      - nav2_compute_path_through_poses_action_bt_node
      - nav2_follow_path_action_bt_node
      - nav2_back_up_action_bt_node
      - nav2_compute_route_bt_node
      # ... (38 plugins total)
\end{lstlisting}

\textbf{Commit:} \texttt{8af6013} (included with Problem 1 fix)

% -----

\subsection{Problem 6: Sparse Road Graph (Drive-Only Network)}

\subsubsection*{Symptom}

The initial campus road graph had only 52 nodes and 113 edges. The nearest graph node to typical goal positions was over 200 meters away, resulting in long indirect routes.

\begin{lstlisting}[language=bash]
# Nearest node to goal (50, 30):
# Node 12 at (255.6, 25.1) --- 205.7m away
\end{lstlisting}

\subsubsection*{Root Cause}

The graph was generated with \texttt{network\_type='drive'}, which only includes roads legally accessible to motor vehicles. The CPP campus has few through-roads but extensive bike paths and walkways.

\subsubsection*{Fix}

Regenerated the graph with \texttt{network\_type='bike'}, which includes all bike-accessible paths:

\begin{lstlisting}[language=bash]
python3 scripts/generate_graph.py \
  --network-type bike \
  --output config/cpp_campus_graph.geojson
\end{lstlisting}

\begin{table}[h]
\centering
\caption{Graph density comparison}
\begin{tabular}{lrrr}
\toprule
\textbf{Network Type} & \textbf{Nodes} & \textbf{Edges} & \textbf{Nearest to (50,30)} \\
\midrule
\texttt{drive} & 52 & 113 & 205.7 m \\
\texttt{bike} & 551 & 1,392 & 35.4 m \\
\bottomrule
\end{tabular}
\end{table}

The bike network provides $10\times$ the node density and $12\times$ the edge density, with significantly better coverage of campus paths.

% ===== SECTION 3: FINAL CONFIGURATION =====
\section{Final Configuration}

\subsection{Behavior Tree}

The final BT uses \texttt{ComputePathToPose} with SmacPlannerHybrid for freespace planning. The Route Server is active and available but not called from the BT (pending first-mile BT nodes on Humble).

\begin{lstlisting}[language=XML,caption={navigate\_route\_graph.xml --- Final BT}]
<root main_tree_to_execute="MainTree">
  <BehaviorTree ID="MainTree">
    <PipelineSequence name="NavigateWithReplanning">
      <RateController hz="0.5">
        <Fallback>
          <ReactiveSequence>
            <Inverter>
              <GlobalUpdatedGoal/>
            </Inverter>
            <IsPathValid path="{path}"/>
          </ReactiveSequence>
          <ComputePathToPose goal="{goal}" path="{path}"
                             planner_id="GridBased"/>
        </Fallback>
      </RateController>
      <FollowPath path="{path}" controller_id="FollowPath"/>
    </PipelineSequence>
  </BehaviorTree>
</root>
\end{lstlisting}

\subsection{Launch Architecture}

\begin{lstlisting}[language=bash,caption={navigation.launch.py node hierarchy}]
navigation.launch.py
  |-- localization.launch.py
  |     |-- sensors.launch.py (URDF, LiDAR, Camera, IMU)
  |     |-- EKF (robot_localization)
  |     +-- NavSat transform (fixed datum)
  |-- actuator_node (cmd_vel -> Teensy UDP)
  |-- nav2_bringup/navigation_launch.py
  |     |-- controller_server (Regulated Pure Pursuit)
  |     |-- planner_server (SmacPlannerHybrid)
  |     |-- smoother_server
  |     |-- behavior_server
  |     |-- bt_navigator
  |     |-- velocity_smoother
  |     +-- lifecycle_manager_navigation
  |-- route_server (nav2_route, GeoJSON graph)
  +-- lifecycle_manager_route (5s delayed start)
\end{lstlisting}

\subsection{Route Server Test Results}

\begin{table}[h]
\centering
\caption{Route computation results --- goal at (50.0, 30.0) in map frame}
\begin{tabular}{lrr}
\toprule
\textbf{Metric} & \textbf{Drive Graph} & \textbf{Bike Graph} \\
\midrule
Planning time & 4.07 ms & 2.37 ms \\
Path poses & 1,445 & 111 \\
Route nodes & 5 & 3 \\
Route edges & 4 & 2 \\
Path length (approx.) & $\sim$720 m & $\sim$55 m \\
Nearest node to goal & 205.7 m & 35.4 m \\
\bottomrule
\end{tabular}
\end{table}

% ===== SECTION 4: COMMITS =====
\section{Commits}

\begin{table}[h]
\centering
\caption{All commits from this session}
\begin{tabular}{lll}
\toprule
\textbf{Hash} & \textbf{Description} & \textbf{Problem} \\
\midrule
\texttt{8af6013} & Use ComputePathToPose in BT + add plugin list & \#1, \#5 \\
\texttt{27fab04} & Set track\_unknown\_space to false & \#2 \\
\texttt{c783af3} & Increase transform\_tolerance to 1.0s & \#3 \\
\texttt{8e7797d} & Add TimerAction delay for route lifecycle & \#4 \\
\bottomrule
\end{tabular}
\end{table}

% ===== SECTION 5: FILES MODIFIED =====
\section{Files Modified}

\begin{table}[h]
\centering
\caption{All files changed during this session}
\begin{tabular}{lp{8cm}}
\toprule
\textbf{File} & \textbf{Changes} \\
\midrule
\texttt{config/navigate\_route\_graph.xml} & Replaced \texttt{ComputeRoute} with \texttt{ComputePathToPose} \\
\texttt{config/nav2\_params.yaml} & Added \texttt{plugin\_lib\_names}; set \texttt{track\_unknown\_space: false}; set \texttt{transform\_tolerance: 1.0} \\
\texttt{launch/navigation.launch.py} & Added \texttt{TimerAction} import and 5s delay for \texttt{lifecycle\_manager\_route} \\
\texttt{config/cpp\_campus\_graph.geojson} & Regenerated with \texttt{network\_type='bike'} (551 nodes, 1392 edges) \\
\bottomrule
\end{tabular}
\end{table}

% ===== SECTION 6: LESSONS LEARNED =====
\section{Lessons Learned}

\begin{enumerate}
    \item \textbf{Humble vs.\ Jazzy differences matter.} BT plugin auto-discovery, first-mile BT nodes, and BehaviorTree.CPP v3 vs.\ v4 format are all version-specific. Always verify against the installed version (\texttt{dpkg -l ros-humble-behaviortree-cpp-v3}).

    \item \textbf{DDS service discovery is not instant.} Lifecycle managers must wait for their managed nodes to register services. A \texttt{TimerAction} delay is the standard pattern for adding servers to existing Nav2 deployments.

    \item \textbf{Costmap defaults assume a static map.} \texttt{track\_unknown\_space: true} is correct when a pre-built map exists. For mapless outdoor operation, \texttt{false} is required so that unobserved space is treated as free.

    \item \textbf{Graph density directly affects route quality.} The \texttt{drive}-only OSMnx network was too sparse for a campus environment. The \texttt{bike} network provides $10\times$ more coverage with paths the vehicle can physically traverse.

    \item \textbf{Test components in isolation.} Calling the route server directly via \texttt{ros2 action send\_goal} before testing through the BT revealed that the server worked correctly --- the bug was in how the BT consumed the result, not in the server itself.
\end{enumerate}

% ===== SECTION 7: NEXT STEPS =====
\section{Next Steps}

\begin{itemize}[noitemsep]
    \item Test full navigation pipeline outdoors with GPS fix and live sensor data
    \item Implement first-mile/last-mile BT logic for Humble (bridge robot position to nearest graph node via \texttt{ComputePathToPose}, then follow graph route)
    \item Configure NTRIP credentials for RTK corrections
    \item Measure physical sensor mount positions and calibrate GNSS lever arm
    \item Evaluate graph quality on the real campus --- prune unusable paths, add missing connections
\end{itemize}

\end{document}
